\documentclass{article}
\usepackage[UTF8]{ctex}
\usepackage{amsmath,amssymb,amsfonts}
\usepackage{enumitem}
\usepackage{geometry}
\geometry{a4paper, margin=1in}
% === 错题本排版优化（不参与题目标识） ===
\setlength{\parindent}{0pt}
\setlength{\parskip}{0.35em plus 0.10em minus 0.08em}
\setlength{\emergencystretch}{3em}
\setlist{itemsep=0.25em, topsep=0.40em, parsep=0pt}
\allowdisplaybreaks[2]
% === 错题本排版优化结束 ===
\begin{document}
\section*{华东师范大学2026年高等代数真题}
\section*{34．华东师范大学2026年高等代数试卷}
一．填空题．每题5分，共10题．\\

\subsection{题目1}
1．如果实系数多项式 \(f(x)=x^{3}+a x^{2}+9\) 有重根，那么这个重根是 \(\_\_\_\_\)．

\subsection{题目2}
2．设矩阵

\[
A=\left(\begin{array}{cccc}
1 & 2 & 0 & 4 \\
-1 & 0 & 2 & -3 \\
5 & 4 & 1 & 0 \\
-1 & -1 & 3 & 2
\end{array}\right) .
\]

其特征值为 \(\lambda_{1}, \lambda_{2}, \lambda_{3}, \lambda_{4}\) ，则 \(\lambda_{1}^{2}+\lambda_{2}^{2}+\lambda_{3}^{2}+\lambda_{4}^{2}=\) \(\_\_\_\_\)．

\subsection{题目3}
3．已知 \(A \in M_{10}(\mathbb{C}), A\) 的特征多项式为 \(x^{5}(x-1)^{5}, A\) 的极小多项式为 \(x^{3}(x-1)^{2}\) ，则 \(A\) 有 \(\_\_\_\_\)个不同的相似等价类．

\subsection{题目4}
4．一个关于未知数 \(\left(x_{1}, x_{2}, x_{3}, x_{4}, x_{5}\right)^{\mathrm{T}}\) 的齐次线性方程组的解空间是由（5，0，2，2，1）\({ }^{\mathrm{T}}\) 和\((-1,7,1,-6,-3)^{\mathrm{T}}\) 张成的线性子空间，那么在 \(x_{1}, x_{2}, x_{3}, x_{4}, x_{5}\) 中，主变元是 \(\_\_\_\_\)．

\subsection{题目5}
5．设线性空间 \(V\) 和它的三个线性子空间 \(V_{1}, V_{1}, V_{3}\) 满足 \(\operatorname{dim} V=9, \operatorname{dim} V_{1}=\operatorname{dim} V_{2}=\operatorname{dim} V_{3}=4\) ， \(V=V_{1}+V_{2}+V_{3},\left(V_{1}+V_{2}\right) \cap V_{3}=\{0\}\) ．则 \(\operatorname{dim}\left(V_{1} \cap V_{2}\right)=\) \(\_\_\_\_\)．

\subsection{题目6}
6．\(\left|\begin{array}{cccccccc}1 & 1 & 1 & 1 & 1 & 0 & 0 & 0 \\ 1 & 2 & 3 & 4 & 0 & 1 & 0 & 0 \\ 1 & 4 & 9 & 16 & 0 & 0 & 1 & 0 \\ 1 & 8 & 27 & 64 & 0 & 0 & 0 & 1 \\ 1 & 8 & 27 & 64 & 0 & 0 & 0 & 1 \\ 2 & 4 & 6 & 8 & 0 & 1 & 0 & 0 \\ 2 & 4 & 6 & 8 & 0 & 1 & 0 & 0 \\ 2 & 16 & 54 & 128 & 0 & 0 & 0 & 1\end{array}\right|=\)

\subsection{题目7}
7．设 \(V\) 是 \(n\) 维线性空间，\(U, W\) 分别是 \(V\) 的 \(m\) 维和 \(r\) 维线性子空间，设 \(T\) 是所有满足 \(f(V) \subset W\)的 \(V\) 的线性变换所构成的线性空间，则 \(\operatorname{dim} T=\) \(\_\_\_\_\)。（用 \(n, m, r\) 表示）\\
\(\_\_\_\_\)

\subsection{题目8}
8．线性空间 \(\mathbb{C}^{3} \times \mathbb{C}^{3}\) 上的线性变换 \(\mathcal{L}\) 定义如下：对 \(X, Y \in \mathbb{C}^{3}, \mathcal{L}(X, Y)=(X+2 Y, 2 X-Y)\) ，则 \(\mathcal{L}\)的所有特征值为 \(\_\_\_\_\)．（若有重数，需写明重数）

\subsection{题目9}
9．设 \(n \geq 3, x_{1}, x_{2}, \cdots, x_{n}\) 是不定元，\(s_{k}\) 是这 \(n\) 个不定元的 \(k\) 次方的和，\(k=1,2, \cdots\) ，用 \(s_{k}\) 表示下述多项式 \(\sum_{1 \leq i<j \leq n} x_{i} x_{j}\left(x_{i}+x_{j}\right)=\)

\subsection{题目10}
10．设 \(V\) 是2025维实线性空间，\(F: V \times V \rightarrow \mathbb{R}\) 是一个非退化双线性型．已知在 \(V\) 的一个线性子空间 \(W\) 上，\(F\) 限制在 \(W \times W\) 上恒为零．考虑所有可能的 \(F\) 和 \(W, \operatorname{dim} W\) 的最大可能值是

二．解答题。每题20分，共5题。\\

\subsection{题目11}
11．考虑关于 \((x, y, z)\) 的线性方程组

\[
\left\{\begin{array}{l}
x+2 y+3 z=0 \\
2 x+3 y+a z=3 \\
3 x+4 y+5 z=b
\end{array}\right.
\]

其中 \((a, b)\) 为参数。试讨论 \((a, b)\) 为何值时，方程组无解，有唯一解，有无穷多解？并在有解（唯一解或无穷多解）时，求出它的所有解。

\subsection{题目12}
12．设 \(V\) 是区间［－1，1］上所有实值连续函数构成的线性空间，定义

\[
(f, g):=\int_{-1}^{1} f(t) g(t) \mathrm{d} t, \forall f(t), g(t) \in V
\]

设 \(W \subset V\) 是由 \(\left\{1, t, t^{2}\right\}\) 张成的线性子空间．\\
（1）利用 Gram－Schmidt 正交化，由 \(\left\{1, t, t^{2}\right\}\) 构造出 \(W\) 的一组规范正交基 \(\left\{w_{1}, w_{2}, w_{3}\right\}\) 。\\
（2）设 \(f(t)=e^{t} \in V\) ，求 \(f(t)\) 在 \(W\) 上的正交投影 \(f_{W}(t)\) ，并用（1）中的规范正交基 \(\left\{w_{1}, w_{2}, w_{3}\right\}\) 表示 \(f_{W}(t)\) ．

\subsection{题目13}
13．设 \(A, B\) 是 \(n\) 阶复方阵，且 \(A\) 可逆．已知存在正整数 \(i\) ，使得对 \(k \in\{t+1, t+2, \cdots, t+n\}, B\) 与\(A^{k}\) 均可交换、证明：\(B\) 与 \(A\) 可交换．

\subsection{题目14}
14．设 \(V\) 是 \(n \geq 2\) 维欧几里得空间，\(v_{1}, v_{2}, \cdots, v_{n} \in V\) ，令 \(A=\left(a_{i j}\right)_{n \times n}\) ，其中 \(a_{i j}=\left(v_{i}, v_{j}\right)\) 是 \(v_{i}\) 与\(v_{j}\) 的内积．\\
（1）证明：\(A\) 是半正定矩阵．\\
（2）证明：\(A\) 是正定矩阵当且仅当 \(v_{1}, v_{2}, \cdots, v_{n}\) 构成 \(V\) 的一组基。\\
（3）若 \(\left(v_{i}, v_{i}\right)=1\) 对任意的 \(1 \leq i \leq n\) ，且 \(\left(v_{i}, v_{j}\right)=c>0\) 对任意的 \(i \neq j\) 。试求 \(\operatorname{det}(A)\) ，并判断 \(c\)满足什么条件时，\(A\) 是正定矩阵．

\subsection{题目15}
15．设 \(A, B\) 均是 \(n\) 阶实对称矩阵，满足 \(A B+A^{3}=B A+B^{3}\) ，证明：\(A=B\) ．
\end{document}
